% partes/parte2/intro.tex
\part{Proceso de decisión de compra digital}

\chapter*{Introducción a la Parte II}
\addcontentsline{toc}{chapter}{Introducción a la Parte II}

La Parte I trató al consumidor como una fuente de datos. La Parte II lo trata
como alguien que decide. Son dos preguntas distintas. La primera pregunta qué
rastro deja una persona. La segunda, cómo llega a comprar.

El programa oficial llama a esta Parte \enquote{Proceso de toma de decisión
de compra digital}, y su competencia es examinar ese proceso a partir del
recorrido y de las figuras que participan en él. Cuatro capítulos la
desarrollan (del 6 al 9). El sexto describe el proceso clásico de cinco etapas y muestra
dónde lo tuerce el entorno digital: reconocimiento inducido por el
algoritmo, reglas de evaluación que cambian la elección y una interfaz que
puede orientar la decisión con información veraz o forzarla con presión
falsa. El séptimo pasa del proceso al recorrido (\textit{customer journey}).
Discute de dónde viene el modelo AIDA, por qué el embudo lineal describe mal
el comportamiento observado y cómo se construye y se limita un mapa del
recorrido; el laboratorio lo modela como una cadena de Markov. El octavo
identifica a las figuras que intervienen (consumidor, usuario, prosumidor,
influenciador, decisor, comprador, evaluador) y muestra con un laboratorio de
redes que influir no es lo mismo que tener muchos seguidores. El noveno
ofrece los criterios, no el catálogo, para elegir las herramientas con que se
analiza todo lo anterior.

Una advertencia metodológica recorre la Parte. Los modelos que aquí se
estudian son representaciones útiles y estilizadas. Ningún mapa de recorrido
captura el camino real de una persona, y varios de los marcos más citados
(el \textit{consumer decision journey} de una consultora, el \textit{messy
middle} de una plataforma publicitaria) son investigación de mercado de
actores con interés comercial. Se presentan como lo que son. Que un marco sea
popular no lo convierte en evidencia.

La práctica oficial de la Parte, \enquote{Recorrido del consumidor digital},
dura 15 horas y tiene un solo entregable. No son cuatro prácticas
independientes: es una misma práctica que crece capítulo a capítulo sobre un
caso real de compra digital que cada equipo elige en el Capítulo 6. Las
actividades de cada capítulo producen un insumo del entregable, y la sección
\enquote{Prácticas oficiales de la unidad} al final de la Parte lo consolida.

Como en toda la obra, cada capítulo cierra con un recuadro de ética y
responsabilidad con los datos. En esta Parte el tema es concreto: reconstruir
el recorrido de una persona exige juntar sus rastros, y juntarlos tiene un
costo para su privacidad y su autonomía.
